\section{Experimental Setup}
\label{sec:experimental_setup}

\subsection{Dataset Split and Preprocessing}

All experiments use the same repository of $N=200000$ simulated trajectories described in Section \ref{sec:data}. A pseudo-random permutation divides the data into $160000$ training trajectories and $40000$ test trajectories. The latter split is used to control the learning-rate scheduler, select the best checkpoint, and report the final metrics.

The scalar features and time-series channels are standardized as described in Section \ref{sec:methods}. Only statistics of the training split are used, and the same saved statistics are applied during all later evaluations. For the GRU models, each raw trajectory is represented by $100$ evenly spaced samples obtained by cubic-spline interpolation. The noise-free trajectory is used to construct all targets, including the impact point in the endpoint-augmented models.

\subsection{Model Configurations}

Table \ref{tab:model_configurations} summarizes the eight configurations. The two MLPs receive the $24$ engineered features introduced in Section \ref{sec:methods}. The first MLP predicts only the Cartesian initial velocity, whereas the second predicts all $11$ launch and environmental parameters.
The GRU prediction head maps the final hidden state through a $128$-unit linear layer, a ReLU activation, dropout with probability $0.1$, and a final linear output layer. Two-layer GRUs additionally use dropout $0.1$ between recurrent layers.
Configurations denoted by \enquote{derivatives} receive the ten channels
\begin{equation*}
    (t,x,y,z,v_x,v_y,v_z,a_x,a_y,a_z),
\end{equation*}
whereas the remaining GRUs receive only $(t,x,y,z)$.

\begin{table*}[htbp]
    \centering
    \small
    \resizebox{\textwidth}{!}{%
    \begin{tabular}{llllll}
        \hline
        ID & Model & Input & Architecture & Output & Crop augmentation \\
        \hline
        01 & MLP velocity & 24 scalar features & $64,32$ & $v_x,v_y,v_z$ & no \\
        02 & MLP all parameters & 24 scalar features & $128,64,32$ & 11 parameters & no \\
        03 & GRU positions & 4 channels & 1 layer, $h=128$ & 11 parameters & no \\
        04 & GRU derivatives & 10 channels & 2 layers, $h=128$ & 11 parameters & no \\
        05 & GRU cropped & 4 channels & 1 layer, $h=128$ & 11 parameters & $f\in[0.2,1.0]$ \\
        06 & GRU cropped derivatives & 10 channels & 2 layers, $h=128$ & 11 parameters & $f\in[0.2,1.0]$ \\
        07 & GRU cropped + endpoint & 4 channels & 1 layer, $h=128$ & 11 parameters + 2D endpoint & $f\in[0.2,1.0]$ \\
        08 & GRU cropped derivatives + endpoint & 10 channels & 2 layers, $h=128$ & 11 parameters + 2D endpoint & $f\in[0.2,1.0]$ \\
        \hline
    \end{tabular}}
    \caption{Model configurations used in the experiments. The symbol $f$ denotes the observed fraction of a trajectory.}
    \label{tab:model_configurations}
\end{table*}

During crop-augmented training, a single prefix length is sampled uniformly between $20$ and $100$ time points for each mini-batch, and the same length is used for every trajectory in that batch. For configurations 07 and 08, the test loss used during training is evaluated at $f=0.5$.
Their two heads minimize
\begin{equation}
    \mathcal{L}
    =
    \f{MSE}\left(\hat{\vb{\lambda}},\vb{\lambda}\right)
    + 0.5\,
    \f{MSE}\left(\hat{\vb{r}}_{\mathrm{hit}},\vb{r}_{\mathrm{hit}}\right),
    \label{eq:endpoint_loss}
\end{equation}
where both target blocks are standardized separately before the loss is computed.

\subsection{Optimization and Model Selection}

All models are optimized with Adam using a batch size of $256$, an initial learning rate of $10^{-3}$, and mean squared error in standardized target space.\cite{kingma2017adammethodstochasticoptimization}
A \texttt{ReduceLROnPlateau} scheduler multiplies the learning rate by $0.1$ after $10$ epochs without improvement, down to a minimum of $10^{-6}$. Early stopping is configured with a tolerance of $32$ non-improving epochs. The checkpoint with the lowest test loss is retained and used in every subsequent evaluation.\cite{bottou2018optimization}

The data permutation is reproducible because its seed is set explicitly. However, due to processing power constraints, only one trained instance of each configuration is available.The reported differences therefore describe these particular training runs. Uncertainty across initializations is not quantified.

The whole machine learning framework is implemented using PyTorch.\cite{paszke2019pytorch}
Additional numerical calculations are performed using the NumPy and SciPy libraries.\cite{harris2020array}\cite{virtanen2020scipy}
Visualizations and plots are realized in Matplotlib.\cite{hunter2007matplotlib}
